% ECONOMICS_PROBLEM: {"id": "C73-649", "title": "One-to-one sink-equilibrium and replicator-attractor correspondence: false", "jel_code": "C73", "source_id": "OP-0149"}
% !TeX program = xelatex
\documentclass[11pt,a4paper]{article}
\usepackage[margin=23mm]{geometry}
\usepackage{fontspec}
\setmainfont{Latin Modern Roman}
\usepackage{amsmath,amssymb,mathtools}
\usepackage{xcolor,enumitem,booktabs,longtable,array,xurl,hyperref}
\definecolor{muted}{HTML}{53636C}
\hypersetup{colorlinks=true,urlcolor=blue,linkcolor=blue}
\setlength{\parindent}{0pt}
\setlength{\parskip}{5pt}
\newcommand{\R}{\mathbb R}
\newcommand{\E}{\mathbb E}
\newcommand{\N}{\mathbb N}
\newcommand{\ind}{\mathbf 1}
\newcommand{\Var}{\operatorname{Var}}
\newcommand{\Cov}{\operatorname{Cov}}
\newcommand{\supp}{\operatorname{supp}}
\DeclareMathOperator*{\argmin}{arg\,min}
\DeclareMathOperator*{\argmax}{arg\,max}
\newcommand{\entrymeta}[3]{{\sffamily\small\color{muted}\textbf{\texttt{#1}}\quad|\quad #2\par #3\par}}
\newcommand{\Problemref}[1]{\texttt{#1}}
\newcommand{\JELref}[2]{\texttt{#2} (JEL \texttt{#1})}
\begin{document}
\section*{One-to-one sink-equilibrium and replicator-attractor correspondence: false}
\textbf{Problem ID:} \texttt{C73-649}\quad \textbf{JEL:} \texttt{C73}\par
% BEGIN ECONOMICS STATEMENT
\label{problem:OP-0149}
% GAME_THEORY_PROBLEM: {"local_id": "GT-019", "original_number": 19, "kind": "H", "entry_kind": "statement_only", "published": false, "review_required": true, "category": "game_theory"}
\label{game:GT-019}
{\sffamily\small\color{muted}
\textbf{GT-019} | Evolutionary dynamics and response graphs\\
\textbf{H} | Historical conjectures explicitly disproved by the cited source.
\par}
\subsection*{Definitions and mathematical statement}
Let $G$ be a finite normal-form game and $X=\prod_i\Delta(S_i)$. Its preference graph has vertices $S=\prod_iS_i$ and an edge $a\to b$ whenever the profiles differ in exactly one player's action and that player's payoff weakly increases: $u_i(b)\geq u_i(a)$. A sink equilibrium is a sink strongly connected component $H$ of this directed graph, with no edge leaving $H$.
The replicator flow on $X$ solves
\[
 \dot x_{i,a}=x_{i,a}\bigl(u_i(a,x_{-i})-u_i(x)\bigr).
\]
Use the source's minimal-attracting-set convention: an attracting set has a strictly larger neighborhood which is forward invariant and all of whose trajectories approach that set; an attractor is inclusion-minimal among such sets. For the stronger correspondence, define
\[
 \operatorname{content}(H)=\{x\in X:\ \prod_i\operatorname{supp}(x_i)\subseteq H\}.
\]
The historical claims were: (i) every replicator attractor contains exactly one sink equilibrium and every sink equilibrium is contained in an attractor; and, more strongly, (ii) the attractors are exactly the sets $\operatorname{content}(H)$. Both claims are false in general finite games.
\subsection*{Short English statement}
The response graph's sink components do not give a one-to-one description of replicator attractors in arbitrary games.
\subsection*{Variants and scope boundaries}
Weak-improvement edges include ties in both directions; substituting strict edges changes the graph. Attractors, chain components, and individual trajectories' omega-limit sets are different objects. Sufficient structural conditions, such as the source's pseudoconvexity condition, define valid restricted results but do not establish the rejected universal correspondence.
\subsection*{Source relationship and status}
The supplied catalog incorrectly describes exact correspondence as conjectural. Version 4 of [S1] states Conjectures 1.1--1.2 as historical claims and disproves them in Section 2, including the two-player case. This corrected section is retained as a historical audit entry.
\subsection*{References}
\begingroup\footnotesize\raggedright
{\footnotesize\raggedright
\textbf{[S1]} Oliver Biggar and Christos Papadimitriou (2026 version). \emph{Sink equilibria and the attractors of learning in games}. arXiv:2502.07975v4, 4 March 2026. Locator: Conjectures 1.1--1.2; Definition 1.3; Sections 2.1--2.3 for counterexamples, Section 3 for pseudoconvexity. \url{https://arxiv.org/html/2502.07975v4}\par}
\par\endgroup

\subsection*{Common conventions: Source mathematical conventions}

\textbf{Finite games.} For a finite set $A$, $\Delta(A)$ is its probability
simplex. Players choose independent mixed strategies in Nash equilibrium;
correlation is allowed only when explicitly part of the solution concept.
In a game with payoff functions $u_i$ and strategy profile $x$, an additive
$\varepsilon$ Nash equilibrium satisfies
\[
 u_i(x)\geq u_i(x_i',x_{-i})-\varepsilon
 \quad\text{for every player $i$ and every admissible deviation $x_i'$}.
\]
For numerical approximation, payoffs are normalized as locally specified.
An $\varepsilon$ payoff residual is distinct from distance at most
$\varepsilon$ to the set of exact equilibria. Point convergence, convergence
of empirical play, and convergence in distance to an equilibrium set are
also distinct.

\textbf{Computational representations.} Unless stated otherwise, a complexity
question takes rational data with binary encoding; polynomial time is measured
in total bit length. A fixed-accuracy polynomial algorithm is not necessarily
polynomial in $1/\varepsilon$, and neither is necessarily polynomial in
$\log(1/\varepsilon)$. Succinct games and value, demand, or sampling oracles
must declare their interfaces. Exact real solutions require an explicit
representation; an unspecified list of arbitrary real numbers is not a
finite computational output. A claimed algorithmic barrier is meaningful
only for its specified model and reductions.

\textbf{Dynamic games.} Histories, information, observation, and allowable
randomization are part of the model. A stationary strategy depends only on
the current state; a positional strategy in a deterministic graph game
chooses one action at each controlled vertex. Nash equilibrium at an initial
state and equilibrium after every history are different requirements.
An expectation of a pathwise $\liminf$ payoff, a $\liminf$ of expected averages,
a limiting expected average, and a uniform finite-horizon equilibrium are
different criteria. None is silently substituted for another.

\textbf{Allocation and mechanisms.} A complete allocation assigns every item
exactly once unless otherwise stated. Zero-valued goods, negative values,
capacity constraints, feasibility, lotteries, and copies of goods affect
fairness definitions and are specified locally. Dominant-strategy incentive
compatibility, Bayesian incentive compatibility, universal truthfulness,
and truthfulness in expectation impose different inequalities. Whenever
an objective ratio has zero denominator, its convention is stated or those
instances are excluded.

\textbf{Infinite-dimensional models.} Expectations, integrals, stochastic
equations, and optimization objectives require their local regularity,
integrability, filtrations, and admissible domains. A backward--forward
mean-field system is a boundary-value problem; it is not an initial-value
semiflow without an additional construction. Long-time, large-population,
and vanishing-noise limits require an order of limits and a convergence
topology. If a minimum need not be attained, the objective is its infimum
or supremum and the approximation requirement is stated separately.
% END ECONOMICS STATEMENT
\end{document}
